% =====================================================================
% Blueprint de checklist - CSCib (Configuracao do ambiente Simulink)
%
% Isto e' apenas o ESQUELETO (secoes + espaco para preencher).
% O conteudo tecnico de cada item deves escrever tu, a partir da
% leitura da Seccao 3 do guia de laboratorio.
%
% Como usar:
%   1. Copia este ficheiro para um novo projeto no Overleaf.
%   2. Em cada "checklist", substitui as linhas em branco pelos
%      passos concretos (consulta as referencias de seccao indicadas
%      em comentario acima de cada bloco).
%   3. Os campos "Sessao" / "Data" no topo servem para reutilizares o
%      mesmo documento em L2, L3, L4-L7, L8-L12, etc.
% =====================================================================

\documentclass[11pt,a4paper]{article}

\usepackage[utf8]{inputenc}
\usepackage[T1]{fontenc}
\usepackage[a4paper,margin=2.5cm]{geometry}
\usepackage{amssymb}   % \square para as caixas de verificacao
\usepackage{enumitem}
\usepackage{titlesec}

% --- Ambiente de checklist reutilizavel ---
\newlist{checklist}{itemize}{1}
\setlist[checklist]{label=$\square$, leftmargin=1.6em, itemsep=0.5em, topsep=0.3em}

\titleformat{\section}{\normalfont\Large\bfseries}{\thesection}{0.8em}{}
\titleformat{\subsection}{\normalfont\large\bfseries}{\thesubsection}{0.8em}{}

\title{Checklist -- Configura\c{c}\~ao do Ambiente de Simula\c{c}\~ao\\\large CSCib -- Laborat\'orio}
\author{}
\date{}

\begin{document}
\maketitle
\vspace{-2em}

\noindent
\textbf{Sess\~ao:} \rule{3cm}{0.4pt} \hfill
\textbf{Data:} \rule{3cm}{0.4pt} \hfill
\textbf{PC usado:} \rule{3cm}{0.4pt}

\vspace{1.5em}

% ---------------------------------------------------------------------
% Ver Seccao 3.2, passos 1-7 (o que NAO precisa de hardware)
% ---------------------------------------------------------------------
\section*{1. Prepara\c{c}\~ao em casa (antes da sess\~ao)}
\begin{checklist}
	\item
	\item
	\item
	\item
\end{checklist}

% ---------------------------------------------------------------------
% Ver Seccao 3.3, passo 10 (Configuration Parameters: Solver,
% Code Generation, Hardware Implementation)
% ---------------------------------------------------------------------
\section*{2. No PC do laborat\'orio -- Ambiente de simula\c{c}\~ao}
\begin{checklist}
	\item
	\item
	\item
\end{checklist}

% ---------------------------------------------------------------------
% Ver Seccao 3.3 (passos 8-9, D/A) e Seccao 3.4 (passos 14-15, A/D)
% ---------------------------------------------------------------------
\section*{3. No PC do laborat\'orio -- Blocos de entrada/sa\'ida}
\begin{checklist}
	\item
	\item
	\item
	\item
\end{checklist}

% ---------------------------------------------------------------------
% Ver "Read this with attention!" no in\'icio do guia
% ---------------------------------------------------------------------
\section*{4. Antes de correr -- Seguran\c{c}a}
\begin{checklist}
	\item
	\item
	\item
\end{checklist}

% ---------------------------------------------------------------------
% O que fazer depois de parar a simula\c{c}\~ao
% ---------------------------------------------------------------------
\section*{5. Ap\'os a experi\^encia}
\begin{checklist}
	\item
	\item
	\item
\end{checklist}

\vspace{1.5em}
\noindent\textbf{Notas / problemas encontrados nesta sess\~ao:}\\[0.5em]
\rule{\linewidth}{0.4pt}\\[1em]
\rule{\linewidth}{0.4pt}\\[1em]
\rule{\linewidth}{0.4pt}

\end{document}
